\section{Introduction}
\label{sec:intro}

One of the main lessons of the HERA programme on deep inelastic scattering (DIS) is the rapid rise of
the gluon density in the proton when Bjorken $x$ decreases at fixed virtuality
$Q^2$~\cite{Aaron:2009aa,H1:2015ubc}. This rise, anticipated on the basis of the BFKL
equation~\cite{Lipatov:1976zz,Kuraev:1977fs,Balitsky:1978ic}, cannot continue indefinitely: at
sufficiently small $x$ the gluon occupation numbers become of order $1/\as$ and non-linear effects ---
gluon recombination and multiple scattering --- tame the growth, leading to gluon
saturation~\cite{Gribov:1984tu,Mueller:1985wy}. The colour glass condensate (CGC) effective
theory~\cite{McLerran:1993ni,McLerran:1993ka,Iancu:2002xk,Iancu:2003xm,Gelis:2010nm,Kovchegov:2012mbw}
provides a systematic framework for this regime: the small-$x$ gluons of the target are described as a
strong classical colour field, whose correlations evolve with rapidity according to the non-linear
Balitsky--Kovchegov (BK)~\cite{Balitsky:1995ub,Kovchegov:1999yj} and
B--JIMWLK~\cite{JalilianMarian:1997jx,JalilianMarian:1997gr,Kovner:2000pt,Weigert:2000gi,Iancu:2000hn,Iancu:2001ad,Ferreiro:2001qy}
equations, while a dilute projectile scatters off this field eikonally, by picking up Wilson lines.

The physics of saturation will be one of the central topics of the future Electron--Ion
Collider~\cite{Accardi:2012qut,AbdulKhalek:2021gbh}. Among the observables that can be measured there,
multi-particle final states in the photon fragmentation region play a special role, because they give
access to correlators of Wilson lines which are more differential than the dipole $S$-matrix probed
by inclusive DIS. Dijet and dihadron production probe the quadrupole and, in the back-to-back regime,
the Weizs\"{a}cker--Williams gluon
distribution~\cite{Gelis:2002nn,Dominguez:2011wm,Metz:2011wb,Dumitru:2015gaa,Mueller:2013wwa}. Trijet
production is the natural next step: it involves six Wilson lines and it is at the same time the
real-emission part of the next-to-leading order (NLO) corrections to dijet production. Leading-order
results for three-parton production in the CGC with massless quarks have been obtained with several
methods~\cite{Ayala:2016lhd,Ayala:2017rmh,Boussarie:2014lxa,Boussarie:2016ogo,Altinoluk:2020qet} and, in
the light-cone wave function approach that we follow here, in ref.~\cite{Iancu:2022gpw}, where it has
been shown that the soft and the collinear limits of the trijet cross-section reproduce the real parts
of the B--JIMWLK and of the DGLAP evolutions, respectively. The complete NLO calculations of dijet and
dihadron production with massless
quarks~\cite{Caucal:2021ent,Taels:2022tza,Bergabo:2022tcu,Caucal:2022ulg,Bergabo:2023wed} and of the
inclusive DIS structure functions~\cite{Balitsky:2010ze,Beuf:2016wdz,Beuf:2017bpd,Hanninen:2017ddy} are
by now available.

Much less is known when the quarks are massive. This is an important gap, for at least three reasons.
First, charm and bottom quarks contribute a sizeable fraction of the DIS cross-section at small $x$ and
are experimentally tagged with good efficiency; heavy-quark dijets and heavy-flavoured hadron pairs are
therefore among the cleanest channels for studying multi-parton correlations. Second, the quark mass
provides a hard scale which keeps the process perturbative down to the photoproduction limit
$Q^2\to 0$, which is the relevant regime for ultra-peripheral nucleus--nucleus collisions. Third, the
mass regulates the collinear singularity of the gluon emission (the ``dead cone''
effect~\cite{Dokshitzer:1991fd}), so that the pattern of logarithms which is resummed by the evolution
of the fragmentation functions differs from the massless case. On the theory side, quark masses have
been included in the NLO light-cone wave functions of the virtual
photon~\cite{Beuf:2021qqa,Beuf:2021srj,Beuf:2022ndu}, with applications to the structure
functions~\cite{Hanninen:2022gje} and to exclusive heavy vector meson
production~\cite{Mantysaari:2021ryb,Mantysaari:2022kdm}, in the NLO JIMWLK Hamiltonian~\cite{Dai:2022imf}
and in the leading-order heavy-quark pair production in the CGC and in its TMD
limit~\cite{Marquet:2017xwy,Altinoluk:2021ygv}. A fully differential expression for the
quark--antiquark--gluon final state with massive quarks, in a form suitable both for trijet
phenomenology and as the real-emission ingredient of heavy-quark dijet production at NLO, is to our
knowledge not available in the literature.

The purpose of this work is to provide such an expression. We consider the process
\begin{equation}
\gT(q)+A\;\longrightarrow\; q(k_1)+\bar q(k_2)+g(k_3)+X\,,
\label{eq:process}
\end{equation}
for a virtual photon with transverse polarisation and for quarks of mass $\mf$, at leading order in
$\as$ and to all orders in the strong colour field of the target. We use light-cone perturbation theory
(LCPT)~\cite{Kogut:1969xa,Bjorken:1970ah,Lepage:1980fj,Brodsky:1997de} to construct the
quark--antiquark--gluon ($\qqg$) Fock component of the \emph{outgoing state} of the photon, that is, of
the state probed by the detector, which includes the collision with the target shock wave. This is the
same strategy as in refs.~\cite{Iancu:2018hwa,Iancu:2020mos,Iancu:2022gpw}, and it is closely related to
the constructions of the soft-gluon light-cone wave function of a fast hadron in
refs.~\cite{Lublinsky:2016meo,Radhakrishnan:2026}, where the LCWF is used to derive the high-energy
evolution itself. Our conventions are those of ref.~\cite{Iancu:2022gpw}, to which our results reduce
when $\mf\to 0$.

The quark mass modifies the calculation in three ways. (i) Both the photon vertex and the gluon-emission
vertex acquire helicity-flip terms proportional to $\mf$. There are therefore four helicity paths
(no flip, flip at the photon vertex, flip at the gluon vertex, double flip) instead of one. (ii) The
energy denominators involve two different mass scales: the usual mass-shifted photon scale
$\ef^2=\zeta(1-\zeta)Q^2+\mf^2$, and a scale $\of=\xi\mf/(1-\zeta)$ which is associated with the gluon
emission and which is responsible for the dead cone. Because of this, the Fourier transform to
transverse coordinates of the amplitude for a gluon emitted \emph{before} the collision is no longer a
single modified Bessel function. We find that each helicity path is described by a radial function with
a compact one-parameter integral representation, eq.~\eqref{eq:Gab}. (iii) The instantaneous
contributions keep their massless form, up to the replacement of $Q^2$ by
$\Ef^2=Q^2+\mf^2(1/\zeta+1/\eta)$.

Our main results are the following.
\begin{itemize}
\item The $\qqg$ component of the outgoing state of the transverse photon with massive quarks,
eqs.~\eqref{eq:state-qbar} and \eqref{eq:state-q}, for an emission by the antiquark and by the quark,
including the instantaneous terms.
\item The trijet cross-section in transverse coordinate space, eqs.~\eqref{eq:xs-qbar},
\eqref{eq:xs-q} and \eqref{eq:xs-int}, with all helicity and polarisation sums performed,
eqs.~\eqref{eq:H-direct}--\eqref{eq:Hx}, and with the colour structure given at finite $\Nc$, at large
$\Nc$ and in the Gaussian approximation.
\item The coherent (elastic) part of the cross-section in momentum space,
eqs.~\eqref{eq:xs-el-qbar}--\eqref{eq:xs-el-int}.
\item The massless limit (section~\ref{sec:massless}), the soft-gluon limit and its relation to the
B--JIMWLK and BK equations (section~\ref{sec:soft}), and the collinear limit, which yields the massive
quasi-collinear splitting function (section~\ref{sec:coll}).
\end{itemize}
All the intermediate and final formulas have been verified by independent symbolic and numerical checks,
which are summarised in appendix~\ref{app:checks}.

The paper is organised as follows. In section~\ref{sec:setup} we fix the kinematics and collect the
LCPT matrix elements and the energy denominators with massive quarks. In section~\ref{sec:wf} we
construct the $\qqg$ component of the outgoing state. Section~\ref{sec:xs} contains the trijet
cross-section and its colour structure, and section~\ref{sec:mom} its coherent part in momentum space.
In section~\ref{sec:validity} we discuss the domain of validity of our results, i.e.\ which
approximations have been made with respect to full QCD. Section~\ref{sec:massless} is devoted to the
massless limit and section~\ref{sec:limits} to the soft and collinear limits, to the corresponding
evolution equations, and to further applications. We conclude in section~\ref{sec:conc}. The appendices
contain the derivation of the radial functions and a summary of the checks.
